
% ═════════════════════════════════════════════════
\section{Basic Git Commands}
% ═════════════════════════════════════════════════

\begin{frame}[fragile]{\textcolor{brown}{Initializing a Repository}}
\begin{block}{\textcolor{white}{Create a new Git repository in your project folder}}
Run the command below to initialize a new git repository.
\begin{lstlisting}[style=gitcode]
git init
\end{lstlisting}

  \vspace{0.4em}
  After initializing, explore with these commands:

  \begin{center}
    \begin{tabular}{ll}
      \hline
      \textbf{Command} & \textbf{Description} \\
      \hline
      \texttt{git status}  & Current state of working directory \& staging area \\
      \texttt{git log}     & Chronological list of all commits \\
      \texttt{git {-}{-}help}  & Display help for Git commands \\
      \hline
    \end{tabular}
  \end{center}
\end{block}
\end{frame}

% ─────────────────────────────────────────────────
\begin{frame}[fragile]{\textcolor{brown}{The Core Workflow}}
  \begin{block}{Staging \& Committing}
    \begin{description}
      \item[\texttt{git add .}] Moves all changes to the \textbf{staging area} (index) — a draft board between your working directory and permanent history.
      \item[\texttt{git commit -m "..."}] Saves a \textbf{permanent snapshot} of staged changes. Use a meaningful message to describe what changed.
    \end{description}
  \end{block}

\begin{lstlisting}[style=gitcode]
git add .
git commit -m "Describe what changed"
git push
\end{lstlisting}

  % \begin{alertblock}{Best Practice}
\emph{\texttt{\textbf{\textcolor{red}{NB: }}}}Commit often with clear messages. You can always restore your project\\[4pt]
to any snapshot.
  % \end{alertblock}
\end{frame}